\documentclass[11pt]{article}

\usepackage[margin=1in]{geometry}
\usepackage[T1]{fontenc}
\usepackage{lmodern}
\usepackage{hyperref}

\title{Synthesis of Long-Range Dependent Symbol Sequences}
\author{Matthew Roughan and collaborators}
\date{\today}

\begin{document}
\maketitle

\section{Annotated Bibliography}

This document starts the literature foundation for S5.jl, a package for
generating long-range dependent non-numerical symbol sequences. The present
version is an annotated bibliography rather than a full paper draft. Manually
downloaded PDFs live under \texttt{paper/downloads/}; open PDFs fetched during
bibliography work live under \texttt{paper/references/}. Items not yet
downloaded, partial downloads, and alternative matches are tracked in
\texttt{DOWNLOAD.mn}.

\subsection{Long-Range Dependence and Self-Similarity}

\paragraph{Samorodnitsky (2006).}
\emph{Long range dependence.}
This survey is a concise bridge between probability theory and applied long-memory
models. It is useful for S5.jl because it separates several meanings that are
easy to blur in finite simulations: covariance-based long memory, heavy tails,
self-similarity, and weak convergence limits. The distinction supports the
package policy of documenting whether a generator has an exact, asymptotic,
latent, nominal, or empirical LRD claim.

\paragraph{Leland, Taqqu, Willinger, and Wilson (1994).}
\emph{On the self-similar nature of Ethernet traffic.}
This paper is one of the central empirical starting points for modern LRD
network-traffic modelling. It showed that packet traffic could retain burstiness
across aggregation scales, motivating models beyond Poisson and short-memory
Markov assumptions. For S5.jl it supplies the broader methodological lesson:
synthetic test data should reproduce large-scale structure, not just local
statistics.

\paragraph{Pipiras and Taqqu (2017).}
\emph{Long-Range Dependence and Self-Similarity.}
This monograph provides the mathematical background tying together covariance
decay, spectral poles, self-similarity, and the Hurst parameter. It is the main
reference for the parameter relationships used throughout the package, including
the links between spectral exponent, ACF decay exponent, and Hurst parameter.

\paragraph{Beran (1994).}
\emph{Statistics for Long-Memory Processes.}
Beran's text is a core statistical reference for inference and modelling in
long-memory time series. Although S5.jl is not an estimator package, the book
helps frame why estimator benchmarking needs synthetic sequences with known
parameters and known violations of short-memory assumptions.

\paragraph{Taqqu, Teverovsky, and Willinger (1995).}
\emph{Estimators for long-range dependence: an empirical study.}
This empirical comparison of LRD estimators motivates the separation between
S5.jl and a future estimator package. The paper demonstrates that estimator
performance depends heavily on model class, sample size, and diagnostics, which
is why S5.jl focuses first on controllable generators.

\paragraph{Witt and Malamud (2013).}
\emph{Quantification of long-range persistence in geophysical time series:
conventional and benchmark-based improvement techniques.}
This survey is useful even though its main examples are geophysical rather than
symbolic. It emphasizes that finite records, trends, nonstationarities, and
method choices can create misleading apparent persistence. That caution maps
directly to S5.jl's validation policy: finite generated sequences should be
benchmarked against known controls, not treated as proof of asymptotic LRD.

\paragraph{Malamud and Turcotte (1999).}
\emph{Self-affine time series: I. generation and analyses.}
This review is a compact reference for generating and analyzing self-affine
numerical time series. It is relevant to the PB generator family because it
collects practical relationships among spectral slopes, Hurst parameters, and
finite-sample diagnostics that are often used before symbolic quantization.

\subsection{Numerical LRD Synthesis}

\paragraph{Paxson (1997).}
\emph{Fast, approximate synthesis of fractional Gaussian noise for generating
self-similar network traffic.}
This is the direct ancestor of S5.jl's PB1 method. It gives a practical spectral
construction of fGn, trading exact covariance matching for speed. In S5.jl the
numerical fGn sequence is mapped to symbols by rank binning, preserving a simple
path from target Hurst parameter to categorical output.

\paragraph{Dieker (2004).}
\emph{Simulation of fractional Brownian motion.}
Dieker's thesis surveys exact and approximate fBm/fGn simulation methods,
including circulant embedding. It provides the natural next reference for
improving PB1 if exact covariance control becomes more important than speed.

\paragraph{Coeurjolly (2000).}
\emph{Simulation and identification of the fractional Brownian motion: a
bibliographical and comparative study.}
This paper is a methodological survey of fBm simulation and estimation methods.
For S5.jl it is a useful checklist for future PB1/PB3 validation, because it
compares multiple synthesis pathways rather than treating one fGn construction
as canonical.

\paragraph{Wood and Chan (1994).}
\emph{Simulation of stationary Gaussian processes in $[0,1]^d$.}
Wood and Chan provide a classic circulant-embedding construction for exact
stationary Gaussian simulation under appropriate embedding conditions. This is a
key reference if S5.jl later adds an exact-covariance Gaussian backend before
mapping numerical samples to symbols.

\paragraph{Dietrich and Newsam (1997).}
\emph{Fast and exact simulation of stationary Gaussian processes through
circulant embedding of the covariance matrix.}
This work complements Wood and Chan by developing fast exact Gaussian-process
simulation through covariance embedding. It is relevant to PB1/PB2 as a possible
upgrade path when exact covariance control matters more than the current fast
approximate spectral path.

\paragraph{Craigmile (2000, 2003).}
\emph{Simulating a class of stationary Gaussian processes using the Davies-Harte
algorithm, with application to long memory processes.}
Craigmile gives a modern treatment of the Davies--Harte method for stationary
Gaussian long-memory simulation. The local download is the 2000 NRCSE technical
report version; the 2003 Journal of Time Series Analysis article is the later
published version. This is a natural reference for comparing PB1's approximate
fGn construction against an exact or near-exact stationary Gaussian baseline.

\paragraph{Sørbye, Myrvoll-Nilsen, and Rue (2017).}
\emph{An approximate fractional Gaussian noise model with $O(n)$ computational
cost.}
This paper approximates fGn with a weighted sum of independent AR(1) components.
It matters for S5.jl because it suggests a scalable path between a simple
finite-state model and an fGn-like covariance target: a small mixture of
short-memory components can approximate long-memory ACFs well over practical
finite ranges.

\paragraph{Roughan, Veitch, and Abry (2000).}
\emph{Real-time estimation of the parameters of long-range dependence.}
This work is important for the wavelet perspective behind the planned PB3
generator. Wavelet methods are attractive because they align naturally with
multi-scale structure and may provide a bridge from LRD synthesis to local
state-machine control. The local download is an extended version of the
IEEE/ACM Transactions on Networking paper, with a related earlier GLOBECOM paper
also retained as supplemental context.

\paragraph{Bai and Shami (2013).}
\emph{Modeling self-similar traffic for network simulation.}
This paper surveys and evaluates practical self-similar traffic models for
network simulation. It is useful implementation context for S5.jl because it
keeps the focus on reproducible generators whose outputs are used downstream by
other experiments, rather than on estimation alone.

\paragraph{Millán (2021).}
\emph{On the generation of self-similar with long-range dependent traffic using
piecewise affine chaotic one-dimensional maps.}
This recent generator paper revisits chaotic-map mechanisms for LRD traffic. It
is not directly a symbolic generator, but it is relevant to MB-style mechanisms:
simple one-dimensional deterministic or nearly deterministic state updates can
produce tunable apparent LRD over finite scale ranges.

\subsection{Point Processes and Heavy-Tailed Models}

\paragraph{Garrett and Willinger (1994).}
\emph{Analysis, modeling, and generation of self-similar VBR video traffic.}
This paper is a classic example of generating self-similar traffic from
heavy-tailed source behaviour. It motivates MB2's heavy-tailed regime-sojourn
construction, where long periods in a regime create large-scale dependence while
the within-regime Markov chain controls local symbol structure.

\paragraph{Lowen and Teich (1995, 2005).}
\emph{Estimation and simulation of fractal stochastic point processes.}
Lowen and Teich provide the point-process foundation for fractal renewal and
shot-noise constructions. The exact 1995 Fractals article has not yet been
downloaded, but a local partial preview of their later book
\emph{Fractal-Based Point Processes} is available and covers closely related
background. Their work informs MB3's view of each symbol stream as an
independent heavy-tailed arrival process.

\paragraph{Thurner, Lowen, Feurstein, Heneghan, Feichtinger, and Teich (1997).}
\emph{Analysis, synthesis, and estimation of fractal-rate stochastic point
processes.}
This locally downloaded arXiv paper is not a replacement for the exact Lowen and
Teich (1995) article, but it is a closely related simulation and estimation study
from the same point-process research line. It is useful context for finite-sample
bias in simulated fractal-rate point processes and therefore for S5.jl's cautious
validation stance.

\paragraph{Ryu and Lowen (1998).}
\emph{Point process models for self-similar network traffic, with applications.}
This paper develops point-process models that connect heavy-tailed mechanisms to
self-similar traffic. It is a key bridge between network-traffic synthesis and
the symbolic FSS construction used in S5.jl.

\paragraph{Roughan, Yates, and Veitch (1999).}
\emph{The mystery of the missing scales: pitfalls in the use of fractal renewal
processes to simulate LRD processes.}
This paper is a warning label for naive heavy-tailed simulation. It shows that
FRP constructions can fail to express LRD over the intended scale range. S5.jl's
validation scripts and TODO list retain this as an explicit design concern for
FSS and On/Off methods.

\subsection{Symbolic and Categorical Models}

\paragraph{Li, Marr, and Kaneko (1994).}
\emph{Understanding long-range correlations in DNA sequences.}
This early review-style paper is directly relevant to non-numerical LRD because
DNA is naturally categorical. It connects empirical DNA correlations with
expansion-modification mechanisms, making it one of the clearest precedents for
symbolic generators whose memory is not merely inherited from a latent numeric
time series.

\paragraph{Salgado-García and Ugalde (2012).}
\emph{Exact scaling in the expansion-modification system.}
This paper analyzes a binary symbolic mutation-replication model and proves
power-law correlation behavior over a range of mutation probabilities. It is
especially useful for S5.jl because it gives a mathematically cleaner symbolic
generation mechanism than many empirical DNA studies.

\paragraph{Colliva et al. (2014).}
\emph{Ising model description of long-range correlations in DNA sequences.}
This paper models DNA correlations using an Ising-style construction. It is a
useful reminder that symbolic biological sequences can be studied through both
direct categorical encodings and latent energy/state models.

\paragraph{Ebeling and Pöschel (1994).}
\emph{Entropy and long range correlations in literary English.}
This paper studies long-range correlations in text using symbolic letters and
entropy-like quantities. It supports an S5.jl use case beyond numerical ACFs:
controlled symbolic sequences can help test whether entropy-rate, mutual
information, and block-complexity diagnostics recover known long-context
structure.

\paragraph{Ogura, Hanada, Amano, and Kondo (2022).}
\emph{Modeling long-range dynamic correlations of words in written texts with
Hawkes processes.}
This paper models word occurrence signals in written texts as Hawkes processes,
separating words with observable dynamic correlations from words that behave
closer to uncorrelated background processes. It is directly relevant to S5.jl's
symbolic focus because it treats word appearances as binary categorical
time-series events and validates simulated point-process outputs against word
occurrence autocorrelation functions. The MB4 `HawkesSymbol` generator is a
finite-history discrete-time adaptation of this self-exciting point-process
idea, not a fitted text model.

\paragraph{Gal, Chen, and Ghahramani (2015).}
\emph{Latent Gaussian processes for distribution estimation of multivariate
categorical data.}
This work motivates PB2. The central idea is to represent categorical outcomes
through latent Gaussian variables and an argmax transformation. S5.jl adapts that
idea to long-memory latent fGn streams and adds finite-sample offset calibration
for marginal control.

\paragraph{Kumar, Raghu, Sarlos, and Tomkins (2017).}
\emph{Linear additive Markov processes.}
This paper motivates MB1. LAMP models make transition probabilities additive in
past observations, allowing long-range memory to be encoded through a decay
kernel over history. S5.jl implements a finite-depth version with a marginal
innovation term.

\paragraph{Singh, Greenberg, and Klakow (2016).}
\emph{The custom decay language model for long range dependencies.}
This language-modelling paper is closely related to LAMP in spirit: it uses a
custom decay over historical context to model long-range dependencies in text.
It supports the idea that symbolic sequences need memory mechanisms distinct
from ordinary short-context Markov models.

\paragraph{Dębowski (2015).}
\emph{Hilberg exponents: new measures of long memory in the process.}
This paper develops power-law growth rates for mutual information and universal
codes. It is relevant to S5.jl because symbolic LRD need not be visible only as
an autocorrelation of arbitrary numeric encodings; information growth can be a
more natural diagnostic for text-like and categorical processes.

\paragraph{Carpio and Daley (2007).}
\emph{Long-range dependence of Markov chains in discrete time on countable state
space.}
This paper is important because it frames LRD for Markov chains through counting
and recurrence-time behaviour. It helps justify the use of count-process and
return-time definitions for non-numerical sequences where ordinary numerical ACF
methods are not directly applicable.

\paragraph{Zeng, Cui, and Li (2025).}
\emph{Fractal illusions: an experimental study of long-range sentence-length
correlations in randomly generated natural language texts.}
This recent paper is valuable as a cautionary reference. It argues that apparent
sentence-length long-range correlations can arise in randomized text generation,
so symbolic LRD diagnostics need controls that distinguish genuine memory
mechanisms from artifacts of preprocessing, punctuation, and measurement.

\paragraph{Dębowski (2025).}
\emph{From Zipf's law to neural scaling through Heaps' law and Hilberg's
hypothesis.}
This recent short paper connects power laws in symbolic language data to
large-scale neural language-model behavior. For S5.jl it is motivation for
keeping symbolic, information-theoretic, and machine-learning validation paths
visible even though the package itself remains a generator library.

\paragraph{Wieczyński and Dębowski (2025).}
\emph{Long-range dependence in word time series: the cosine correlation of
embeddings.}
This recent title is relevant to categorical sequence analysis because it
studies word sequences through embedding-based correlations rather than direct
numeric observations. It suggests a possible validation direction for symbolic
generators whose alphabet elements later receive semantic or learned
representations.

\subsection{Entropy, Predictability, and Information Storage}

\paragraph{Crutchfield and Feldman (2003).}
\emph{Regularities unseen, randomness observed: Levels of entropy convergence.}
This paper is a core reference for excess entropy and entropy-rate convergence.
It frames excess entropy as stored predictive information in a process, making it
a natural quantity to test on controlled LRD symbolic sequences. S5.jl can create
cases where local Markov structure and long-range structure are varied
separately, which is exactly the kind of setting needed to probe entropy-rate
convergence.

\paragraph{Bialek, Nemenman, and Tishby (2001).}
\emph{Predictability, complexity, and learning.}
This work connects prediction, complexity, and information in sequences. It is
useful background for studying how much long-range structure a learner can
extract from finite data, and for interpreting excess entropy-like quantities as
predictive information.

\subsection{Machine Learning Motivation}

\paragraph{Belletti, Beutel, Jain, and Chi (2018).}
\emph{Factorized recurrent neural architectures for longer range dependence.}
This paper motivates the project from the machine-learning side by showing that
common recurrent architectures struggle with long-range dependence. It helps
explain why synthetic NN-LRD data is useful for stress-testing learning
algorithms.

\paragraph{Belletti, Chen, and Chi (2019).}
\emph{Quantifying long range dependence in language and user behavior to improve
RNNs.}
This paper is one of the direct motivations for studying LRD in language and
human behavioural sequences. It connects empirical LRD measurement to model
design, reinforcing S5.jl's role as a controllable data source for later
estimator and learning experiments.

\paragraph{Vaswani et al. (2017).}
\emph{Attention is all you need.}
The transformer architecture made long-context neural sequence modelling central
to modern machine learning. Although S5.jl is not language-specific, its symbolic
generators can create non-language sequences with controlled long-context
dependencies for training and stress-testing transformer-style models.

\paragraph{Tay et al. (2021).}
\emph{Long Range Arena: A benchmark for efficient transformers.}
Long Range Arena motivates synthetic and semi-synthetic tasks designed to test
long-context modelling. S5.jl extends this spirit to LRD categorical processes
where statistical structure, alphabet, marginal distribution, and local Markov
structure can be controlled.

\subsection{Two-Dimensional and Spatial Self-Similar Processes}

\paragraph{Gelbaum (2012).}
\emph{Fractional Brownian fields over manifolds.}
This work generalizes fractional Brownian fields beyond one-dimensional time.
It is outside S5.jl's current scope, but it is a useful reference for future
thinking about spatial categorical fields or image-like symbolic processes where
LRD appears across two or more dimensions.

\paragraph{Gelbaum and Titus (2013).}
\emph{Simulation of fractional Brownian surfaces via spectral synthesis on
manifolds.}
This short paper provides a spectral-synthesis construction for random fractal
surfaces. It belongs in a 2D reference section because it shows how the spectral
ideas behind PB1 can be generalized to surfaces while preserving an explicit
simulation pathway.

\paragraph{Balcerek et al. (2025).}
\emph{Two-dimensional fractional Brownian motion: analysis in time and frequency
domains.}
This recent paper studies two-dimensional fBm with dependent components,
including covariance and spectral structure. It is not a target for S5.jl's
current 1D symbolic API, but it is useful context if the project later considers
multi-dimensional symbolic processes or vector-valued latent drivers.

\paragraph{Stein (2002).}
\emph{Fast and exact simulation of fractional Brownian surfaces.}
Stein's paper is a classic reference for exact simulation of fractional Brownian
surfaces. It is a useful benchmark reference for any future extension from
one-dimensional symbolic sequences to spatial or image-like symbolic fields.

\end{document}
